Our entire analysis will be conditioned on the occurrence of several events,
each of which occurs with probability approaching one as $N$ goes to infinity.
Outside of these events we provide no guarantees.  As a consequence, the
residuals of our approximation can be used to show convergence in probability
or distribution (via Slutsky's theorem), but not convergence in expectation
nor almost sure convergence.

% \begin{defn}\deflabel{regular_event}
%

% \end{defn}

%%%%%%%%%%%%%%%%%%%%%%%%%%%%%%%%%%%%%%%%%%%
%%%%%%%%%%%%%%%%%%%%%%%%%%%%%%%%%%%%%%%%%%%

\begin{lem}\lemlabel{regular_event}
    %
For any $\epsilon_U > 0$, let $\regevent$
denote the event that all of the following occur
for a particular $N$, $\delta_1$, $\delta_2$.
(``R'' is for ``regular''.)
%
\begin{enumerate}
%
\item \eventlabel{disjoint_regions} The regions $R_1$, $R_2$, and $R_3$ are
disjoint, and $I(\phi) = I_1(\phi) + I_2(\phi) + I_3(\phi)$.
%
\item \eventlabel{thetahat_close}
$\norm{\thetahat - \thetatrue}_\infty < \epsilon_U$.
%
\item \eventlabel{prior_close}
    $|\prior(\thetahat) - \prior(\thetatrue)| < \epsilon_U$.
%
\item \eventlabel{ulln} For any scalar-valued function $\rho(\theta, \x_n)$
with $ \expect{\fdist(\xn)}{\rho(\theta, \xn)} = 0$
satisfying a ULLN with $\delta = \deltaulln$, we have
%
%\begin{align*}
%
$\sup_{\theta \in R_1 \bigcup R_2} \meann
\abs{\rho(\theta, \x_n)}
    \le \epsilon_U$.
%
\item \eventlabel{infoev} $\infoevhat > \infoev / 2$.
%
\item \eventlabel{strict_opt} There exists $\epsilon_L > 0$ such that
%
%\begin{align*}
$\sup_{\theta \in R_3}
\left( \sumn \ell(\x_n \vert \thetahat) -
       \sumn \ell(\x_n \vert \theta)
\right) > N \epsilon_L$.
%\end{align*}
\end{enumerate}

Let \assuref{bayes_clt} hold. For any $\delta_1 > 0$, any $\pi_0 \in (0,1)$, and
any $\delta_2 > 0$ and $\epsilon_U > 0$ satisfying  $\deltaulln > \epsilon_U +
\delta_2$ and $\delta_2 / 2 \ge \epsilon_U$, there exists an $N^*$ such $N >
N^*$ implies that event $\regevent$ occurs with $\fdist$-probability at least $1 -
\pi_0$.
%
\begin{proof}
%
\Eventref{disjoint_regions} follows from the fact that, for any positive
$\delta_1$ and $\delta_2$, eventually $\delta_2 \sqrt{N} > \delta_1 \log N$.

\Eventref{thetahat_close} follows from \lemref{thetahat_consistent}.

\Eventref{prior_close} follows from the fact that the prior density is
continuous by \assuref{bayes_clt} \itemref{prior_smooth} and,
by \lemref{thetahat_consistent}, we can take
$\thetahat$ to be arbitrarily close to $\thetatrue$.

For \eventref{ulln}, observe that $\norm{\thetahat - \thetatrue}_2 < \epsilon_U$
implies that
%
\begin{align*}
%
\theta \in& R_1 \bigcup R_2 \Rightarrow\\
\norm{\theta - \thetahat}_2 \le& \delta_2 \Rightarrow\\
\norm{\theta - \thetatrue}_2 \le&
    \norm{\theta - \thetahat}_2 + \norm{\thetahat - \thetatrue}_2
\le \delta_2 + \epsilon_U \le \deltaulln.
%
\end{align*}
%
Therefore,
%
\begin{align*}
%
\sup_{\theta \in R_1 \bigcup R_2} \meann \abs{\rho(\theta, \x_n)} \le&
\sup_{\theta \in \thetaball{\deltaulln}}
    \meann \abs{\rho(\theta, \x_n)} \plim 0.
%
\end{align*}
%
\Eventref{infoev} follows from \lemref{thetahat_consistent}.


For \eventref{strict_opt}, write
%
\begin{align*}
%
\MoveEqLeft
\frac{1}{N} \sup_{\theta \in R_3}
\left( \sumn \ell(\x_n \vert \thetahat) -
       \sumn \ell(\x_n \vert \theta)
\right)
=\\&
\likhat(\thetahat) - \logprior(\thetahat) -
\sup_{\theta \in R_3} \left(\likhat(\theta) - \logprior(\theta)\right)
=\\&
\likhat(\thetahat) - \likhat(\thetatrue) +
    \sup_{\theta \in R_3}
        \left(\likhat(\thetatrue) - \likhat(\theta)\right).
%
\end{align*}
%
Since $\norm{\thetahat - \thetatrue}_2 \le \epsilon_U \le \delta_2 / 2$
by assumption of the present lemma,
%
\begin{align*}
%
\theta \in& R_3 \Rightarrow\\
\norm{\theta - \thetahat + \thetahat - \thetatrue}_2 \ge & \delta_2 \Rightarrow\\
\norm{\theta - \thetahat}_2 + \norm{\thetahat - \thetatrue}_2
    \ge& \delta_2 \Rightarrow\\
\norm{\theta - \thetahat}_2 \ge& \delta_2 - \frac{\delta_2}{2} = \frac{\delta_2}{2},
%
\end{align*}
%
so the set $R_3$, is contained in $\thetadom \setminus \thetaball{\delta_2 /
2}$. Consequently, by \assuref{bayes_clt} \itemref{strict_opt}, there exists an
$\epsilon_L$ such that
%
\begin{align*}
%
\sup_{\theta: \norm{\theta - \thetahat}_2 \ge \delta_2}
    \left(\likhat(\thetatrue) - \likhat(\theta)\right) \le&
\sup_{\theta \in \thetadom \setminus \thetaball{\delta_2 / 2}}
    \left(\likhat(\thetatrue) - \likhat(\theta)\right) \le
-2 \epsilon_L.
%
\end{align*}
%
By the continuity of $\lik$, and a ULLN applied to $\likhat(\theta)$, we can
take $\norm{\thetahat - \thetatrue}_2$ sufficiently small that that
$|\likhat(\thetahat) - \likhat(\thetatrue)| < \epsilon_L$.  It follows that,
with probability approaching one,
%
\begin{align*}
%
\frac{1}{N} \sup_{\theta: \norm{\theta - \thetahat}_2 \ge \delta_2}
\left( \sumn \ell(\x_n \vert \thetahat) -
       \sumn \ell(\x_n \vert \theta)
\right) \le \epsilon_L - 2 \epsilon_L = -\epsilon_L.
%
\end{align*}
%
The result follows by multiplying both sides of the inequality by $N$.
%
\end{proof}
%
\end{lem}

Note that \lemref{regular_event} requires $\delta_2$ to be sufficiently small (so that
$\deltaulln > \delta_2$), which in turn requires $\epsilon_U$ to be small (so
that $\deltaulln > 3 \epsilon_U$).  For any particular $\deltaulln > \delta_2$,
we can always apply \lemref{regular_event} with arbitrarily small $\epsilon_U$, at the
cost of requiring a large enough $N^*$ so that the probability convergence
results apply.
%
Throughout, we will choose $\epsilon_U$ \textit{small enough} that we obtain our
desired results.
